\chapter{Análisis y Metodología}
\label{sec:analisis_metodologia}
En esta sección se presenta el marco metodológico empleado para analizar la estructura de dependencia entre activos financieros mediante técnicas de Hierarchical Risk Parity (HRP) y su extensión cuántica. El objetivo principal es estudiar las relaciones latentes del mercado utilizando métricas alternativas a la correlación tradicional, permitiendo una representación más rica y robusta de las interdependencias entre activos.

\section{Representación multivariante de los activos}
\label{sec:representacion_multivariante}

Cada activo financiero se representa mediante un conjunto de características económicas que describen distintos aspectos de su comportamiento dinámico. Para cada activo $i$ se construye una matriz de observaciones:

\[
\mathbf{X}_i \in \mathbb{R}^{T \times P},
\]

donde $T$ denota el número de observaciones temporales y $P$ el número de características consideradas. En este trabajo se emplean seis características por activo, diseñadas para capturar información relacionada con retornos, volatilidad, momentum, reversión y volumen negociado.

El uso de una representación multivariante permite superar las limitaciones del enfoque clásico basado exclusivamente en la correlación de retornos, la cual es conocida por su elevada inestabilidad temporal y sensibilidad al ruido.

\section{Normalización de características}
\label{sec:normalizacion_caracteristicas}
Con el fin de garantizar una contribución equilibrada de todas las variables, las características se estandarizan de forma independiente para cada activo. Este procedimiento evita que determinadas magnitudes, como la volatilidad o el volumen, dominen el proceso de modelado y asegura la comparabilidad entre activos.


\section{Representación cuántica mediante matrices densidad}
\label{sec:representacion_cuantica}
Para modelar relaciones no lineales entre activos, cada observación multivariante se embebe en un espacio cuántico mediante un mapeo de características cuántico. A partir de dicho mapeo, cada observación queda representada como un estado cuántico puro:

\[
\ket{\psi(\mathbf{x}_i^t)}.
\]

Posteriormente, para cada activo se construye una matriz densidad promedio definida como:

\[
\rho_i = \frac{1}{T} \sum_{t=1}^{T}
\ket{\psi(\mathbf{x}_i^t)}\bra{\psi(\mathbf{x}_i^t)}.
\]

Esta formulación permite capturar la distribución completa de las observaciones del activo en el espacio de Hilbert, proporcionando una representación más informativa que los estadísticos clásicos de segundo orden.

\section{Métrica de distancia cuántica}
\label{sec:metrica_distancia_cuantica}
La similitud entre activos se cuantifica mediante la distancia de Frobenius entre sus matrices densidad promedio. Para dos activos $i$ y $j$, dicha distancia se define como:

\[
d_F(\rho_i, \rho_j) = \frac{1}{2}
\sqrt{\mathrm{Tr}\left[(\rho_i - \rho_j)^2\right]}.
\]

Esta métrica proporciona una medida geométrica natural en el espacio de operadores cuánticos y permite comparar activos teniendo en cuenta la totalidad de su comportamiento multivariante.

\section{Construcción y ordenamiento jerárquico}
\label{sec:construccion_ordenamiento}

Las distancias cuánticas entre todos los pares de activos se organizan en una matriz de distancias simétrica. Dicha matriz se emplea como entrada para un procedimiento de clustering jerárquico basado en distancia mínima, a partir del cual se obtiene un dendrograma que describe la estructura de agrupación del mercado.

El orden de los activos se extrae de las hojas del dendrograma y se utiliza para reorganizar las matrices de distancia. Este reordenamiento no altera los valores originales, pero permite revelar patrones estructurales y agrupaciones coherentes entre activos.

\section{Visualización de la estructura del mercado}
\label{sec:visualizacion_estructura}

Para analizar visualmente la estructura resultante, se representan las matrices de distancia cuántica antes y después del ordenamiento jerárquico mediante mapas de calor.

La Figura 3.1 muestra la matriz de distancia cuántica sin ordenar, en la cual no se observa una estructura aparente debido a la disposición arbitraria de los activos. Por el contrario, la Figura 3.2 presenta la misma matriz tras el ordenamiento jerárquico, donde emergen agrupaciones coherentes que reflejan similitudes estructurales entre activos.

Con el objetivo de mejorar la interpretabilidad visual, se emplea un escalado robusto de la escala de colores basado en percentiles, el cual reduce la influencia de valores extremos sin modificar la información subyacente.

\begin{figure}[Ejemplo de uso de figure]{FIG:EJEMPLO}{Matriz de distancia cuántica sin ordenar. La disposición arbitraria de los activos impide identificar patrones estructurales.}
        \image{13cm}{}{fig3-corr-matrices}
\end{figure}

\begin{figure}[Ejemplo de uso de figure]{FIG:EJEMPLO}{Matriz de distancia cuántica tras el ordenamiento jerárquico. El reordenamiento revela agrupaciones coherentes de activos, evidenciando la estructura latente del mercado capturada por la representación cuántica.}
        \image{13cm}{}{fig4-quantum-distance}
\end{figure}

El contraste entre ambas representaciones permite apreciar el valor del ordenamiento jerárquico como herramienta de análisis estructural, facilitando la interpretación cualitativa de las relaciones complejas entre activos financieros.
